\part{Unit 4 - Clustering}
\chapter{Introduction}
In this Unit, we will  investigate so-called mixture models. In mixture models
it is assumed that a data point can be generated by any of a number of
stochastic processes. An often used example is a mixture of Gaussians, but mixtures of other models are
also possible, e.g. mixtures of Bernoulli processes.

\begin{figure}[!ht]
  \begin{center}
    \includegraphics[width=0.4\textwidth]{mixtureunlabeled.pdf}
    \caption{An example of a dataset for which a mixture of Gaussians can be considered.}
    \label{fig-mix}
   \end{center}
\end{figure}


These mixtures are very important because they allow us to model complex multimodal distributions, which we often
cannot handle analytically in terms of a collection of Gaussian processes for which we have familiar tools available.
Here, an interesting problem arises: if we have a complex distribution modelled as a mixture of Gaussians and we
have a dataset that looks like it has a set of clusters (have a look at Fig. \ref{fig-mix} as an example), then
we are faced with a double problem. First, we must infer the parameters of the Gaussians that are supposed to make
up the cluster. We know how to do that for a single Gaussian since we can do maximum likelihood estimation. But here we
also must somehow infer how much each individual Gaussian contributes to the data cloud, or more precisely: we must infer
the probability that a given data point comes from a certain cluster, as well as estimating the cluster properties themselves,
and ideally, do this simultaneously. In order to even start that process,  we have to make a guess about $K$, the number of clusters.

In Fig. \ref{fig-mix} it looks like $K=3$ is a reasonable choice, but in a data space with more than two dimensions we may struggle
to visualise the complete dataset and the number of clusters may not be readily estimated. This is an important problem,
which we shelf for now. We assume that we have determined the number of clusters, $K$, that we want to use to model our data,
which lives in a space of dimension $D$, $D=2$ in this case.

We can represent the probability that a given data point $\boldsymbol{x}$ belongs to cluster $j$ by a $K$-dimensional
vector $(\pi_1, \cdots, \pi_K)^T$. So the probability that a given data point belongs to a certain cluster is a \emph{stochastic variable}.
The problem is that we are not observing that variable: we are not given a label that tells which cluster this data point belongs to. A stochastic
variable whose realisation cannot be observed but has to be inferred through indirect means is called a \emph{latent variable}. What
we will encounter in the Unit is probability distribution functions that are written as $(\boldsymbol{X}, \boldsymbol{Z})$, which
is a joint probability over variables $X$, whose outcome can be observed directly and $Z$, whose outcomes cannot be observed.

Inferring this joint distribution is clearly much harder to do than the kind of inference in Unit 1, where we had to estimate the parameters
of known distributions, but where we were able to produce estimators, either maximum likelihood estimators or posterior distributions
that made direct use of the data observations. Here, we have to proceed in an iterative process that consists of two steps:
the first step is to make a good guess about the distribution of the latent variables. Then, having numerical values for the latent variable, we
can proceed with conventional methods for estimating the observable variables. We then repeat and try to use these estimates
to make better estimates for the latent variables, and so on. This leads to a class of algorithms called \emph{E/M algorithms}. We will
investigate this algorithm in detail for the \emph{Gaussian Mixture Model} (GMM).

The presence of latent variables is not just a problem. Some of the recent progress in AI and machine learning is a new approach
to look at complex high dimensional data, like images or large volumes of text, as being 'caused' by simpler underlying
processes. These methods try to use a large number of observations to estimate the joint probability $p(X,Z)$ where
$X$ is high dimensional (think images) and $Z$ has much smaller dimensionality. The aim becomes to learn $p(Z | X)$. In practice, this is so hard
that one tries to create a model distribution $q(Z \mid X)$ that is a good approximation.  If this process is successful, then 
one has access to the latent distribution $p(Z)$. This process is called \emph{encoding} and we will see that GMMs are
a very rudimentary form of this. Encoding is useful in its own right. It has been used successfully in \emph{image compression}.

Modern techniques go one step further.
Learning the mapping $p(X \mid Z)$ is called \emph{decoding}. If one successfully learns this mapping,
it becomes possible to generate synthetic images like that resemble ones from the original dataset
from samples that have been generated in the latent space.
The theoretical framework underpinning all of these methods is called \emph{variational inference}. In variational inference the objective
is to approximate the true latent variable distribution $p(Z \mid X)$, which is often intractable  by an estimate $q(Z)$ which
has a simpler functional form.
To monitor whether this training is successful, we need something called the \emph{Evidence Lower Bound} (ELBO).

The purpose of this chapter is to provide:
\begin{enumerate}
\item The $K$-means clustering algorithm
\item An introduction to GMMs and E/M algorithms.
\item An introduction to the concepts ELBO and \emph{posterior latent variable distribution}.
\item A brief overview of where in modern machine learning they feature.
\item Pointers to the literature where more technical expositions of the variational framework are available.
\end{enumerate}
A full discussion of the variational framework is well outside the scope of this module but a gentle introduction to some of its concepts is useful for
modern machine learning.



\chapter{K-Means}
\section{K-Means: a Simple Clustering Algorithm}
In the \emph{Data Science} module the K-Means algorithm was introduced. We will discuss it here briefly to
set the stage of Gaussian mixture models. Once we have discuss those, we will see that K-Means is a special case.

The K-Means algorithm assumes that  a dataset can be partitioned into $K$ clusters, where $K$ is a number that
has to be determined by the data analyst. Once the algorithm is run, each data point has been assigned to a cluster
and a mean position for each cluster has been calculated.

The algorithm is straightforward. Assuming there are $N$ data points $\boldsymbol{x}_i$ in a $D$-dimensional space, a one-over-$K$ coding will indicate to which of the clusters
the point belongs. This is recorded in a matrix a $N \times K$, $\boldsymbol{R}$ whose components are $r_{ik}$, where $i$, the row index, labels the
data point and $k$ the cluster. For a given $i$, all but one of the $r_{ik}$ are zero, the non-zero component is equal to 1. You can think of this
as large matrix with as many rows as there are data points where each row contains a one-over-K coding labelling which cluster the data point belongs to.

The K-means algorithm aims to minimise the \emph{distortion measure}:
\begin{equation}
  J = \sum^N_{i=1} \sum^K_{k=1} r_{ik} \mid \mid \boldsymbol{x}_i  - \boldsymbol{\mu}_k \mid \mid^2
\end{equation}

This represents the sum of squares of distances from each data point to some cluster centre $\boldsymbol{\mu}_k$. Initially, we know neither
cluster centres nor cluster assignment. We start out by simply picking random values for the cluster values $\boldsymbol{\mu}_k$. Then we repeat the following
two steps until hitting some stopping criterion:

\begin{enumerate}
\item Assign each data point to a clusters. This is done by pick $\boldsymbol{R}$ as follows:
$$
r_{ik} = \left\{ \begin{array}{cc} 1 & \mbox{if } k = \mbox{arg min}_j  \mid \mid \boldsymbol{x}_i -\boldsymbol{\mu}_j \mid \mid^2 \\ 0 & \mbox{otherwise} \end{array} \right.
$$
\item Estimate the new cluster centres as follows:

$$
\boldsymbol{\mu}_k = \frac{\sum_i r_{ik} \boldsymbol{x}_i}{\sum_i r_{ik}}
$$
\end{enumerate}

A stopping criterion can be defined by demanding  that $J$ be below a certain value, or by visual inspection of the resulting clusters when possible.

\begin{figure}[!ht]
  \begin{center}
    \includegraphics[width=0.7\textwidth]{of.pdf}
  \end{center}
  \caption{The K-Means algorithm on the Old Faithful dataset. Despite the unfortunate choice for the starting values of the centres, the algorithm converges
    in approximately 5 steps, leaving two well separated clusters. The black marker represents the centre estimate for each cluster at the start
  of the algorithm. The red crosses the new position at the end of each step.}
  \label{fig-of}
\end{figure}

In Figure \ref{fig-of} we show the convergence of the algorithm on the so-called \emph{'Old-Faithful'} dataset.

On separated clusters this algorithm works fairly well. Importantly, it already displays the two-step approach that characterises E/M algorithms, as we
will see later.

One disadvantage of this algorithm is that we need a good idea about how many clusters there are. If we assume the wrong number we can create
extra clusters that are not really separate. Another disadvantage is that when we know that clusters overlap this algorithm cannot respect the overlap
because points are assigned to some cluster based on a notion of distance to some centre and therefore partition boundaries. This can happen
when different processes can generate data samples that overlap in data space. The K-means algorithm has to make a hard assignment of a point
to each cluster, which can deform the cluster and affect its estimates (see \emph{Activity 5.a}).

\begin{figure}[!ht]
  \begin{center}
    \includegraphics[width=0.4\textwidth]{mixtureunlabeled.pdf}
    \includegraphics[width=0.4\textwidth]{mixturelabeled.pdf}
    \caption{The plot on the left shows a situation that occurs often in practice. The data appears to be clustered but without any
      labels it is not straightforward to assign points to clusters and to calculate the cluster properties. When we do show
      the labels, as in he right plot, it is clear that different clusters can inhabit the same part of data space, which the standard $K$-means clustering
      does not recognise.}
    \label{fig-mixtexample}
   \end{center}
\end{figure}


\chapter{Mixtures of Gaussians}

A mixture of Gaussian provides a good example of clusters that inhabit the same part of data spaces.
In a mixture of Gaussians, a data point can be created by any number of Gaussian distributions. To each
Gaussian, we can assign a probability that indicates how likely a particular Gaussian generates the next data point. Out of these Gaussian a single one is selected,
which is then sampled subsequently to generate the actual data point. Figure \ref{fig-mixtexample} shows and example of a dataset that has been generated this way.

In order to model such a data generation process, we need a generalisation of the Bernoulli process introduced in Unit 1. We assume that there is a $K$-dimensional
vector $\boldsymbol{x}$, where one of the elements $x_{i}$ has the value one and the other elements are 0. We have already encountered this as the 1-over-$K$ coding
scheme. In order to model it, we need a multidimensional version of the Bernoulli process. We define a set of $K$ probabilities, $\mu_k$, which by virtue of them probabilities
must add to one:
$$
\sum_k \mu_k = 1
$$

The probability to generate a particular 1-over-$K$ pattern $\boldsymbol{x}$ is given by:
\begin{equation}
  p(\boldsymbol{x} \mid \boldsymbol{\mu} ) = \Pi^K_{k=1} \mu^{x_k}_k
  \label{eq-berngen}
\end{equation}

The way to read this is that if we want to know the probability, given $K=5$, what the probability for $(0,1,0,0,0)$ to occur is
that the only $x_k \ne 0$ is $x_2$ = 1. The probability for this particular pattern is therefore $\mu_2$, in other words Eq. \ref{eq-berngen}
is a succinct mathematical expression of what we just explained in words. Note that because one and only one of the $x_k =1$, rather trivially it follows
that $\sum^K_{k=1} x_k = 1$.

This distribution is properly normalised:
$$
\sum_{\boldsymbol{x}} p( \boldsymbol{x} \mid \boldsymbol{\mu} ) = \sum^K_{k=1} \mu_k = 1
$$
and that the expectation value is given by:
$$
\mathbb{E}[ \boldsymbol{x} \mid \boldsymbol{\mu} ] = \sum_{\boldsymbol{x}} p(\boldsymbol{x} \mid \boldsymbol{\mu}) \boldsymbol{x} = (\mu_1, \cdots \mu_K)^T = \boldsymbol{\mu}
$$

As with the Bernoulli process, we can write down  likelihood for a dataset $\mathcal{D} = \boldsymbol{x}_1, \cdots \boldsymbol{x}_N$, where the $\boldsymbol{x}_i$ indicate
$N$ independent observations:
$$
p(  \mathcal{D} \mid \boldsymbol{\mu} ) = \Pi^N_{i=1} \Pi^K_{k=1} \mu^{x_{ik}}_k = \Pi^K_{k=1} \mu^{\sum_i x_{ik}}_k = \Pi^K_{k=1} \mu^{m_k}_k
$$
Here
$$
m_k = \sum_i x_{ik},
$$
i.e. the number of observations of $x_k = 1$.

We can find the maximum likelihood estimator (MLE) for $\boldsymbol{\mu}$ by maximising this expression, subject to the condition that $\sum_k \mu_k = 1$. You may
or may not be familiar with the technique of Lagrangian multipliers, which is the appropriate technique for finding constrained optima. If you have not seen this
you can just take the result at face value; we will not assess you on this. If you are curious, a brief exposition is given in appendix E of \cite{bishop2006}.

The technique consists of multiplying the constraint by an arbitrary factor, the so-called multiplier and then adding to the quantity whose extremum we want to find.
Again, we work with the log likelihood which means that we want to maximise the quantity:
$$
\sum^K_{k=1} m_k \ln \mu_k,
$$
subject to $\sum^K_{k=1} \mu_k = 1$.
Adding the constraint term we get:
$$
\sum^K_{k=1} m_k \ln \mu_k + \lambda( \sum^K_{k=1} \mu_k - 1)
$$
Optimising with respect to $\mu_j$ gives:
$$
\mu_j = -\frac{m_j}{\lambda}
$$
Substituting this into the constraint that all $\mu_j$ must add to one,
we find that:
$$
- \sum_k \frac{m_k}{\lambda} = 1, 
$$
or $\lambda = -N$ so that ultimately we find the common sense result:
$$
\mu_k = \frac{m_k}{N}
$$

In passing we mention that for the multidimensional Bernoulli distribution a generalised version of the Beta distribution exists, which is called the \emph{Dirichlet} distribution.
One would use this distribution as a prior for Bayesian inference, see section 2.2.1 of \cite{bishop2006} for details.




We are now ready to present the GMM model:
We can interpret this as a generative model: first we use a Bernoulli process to select a cluster $i$, then we sample from the Gaussian with mean $\boldsymbol{\mu}_i$,
covariance matrix $\boldsymbol{\Sigma}_i$.

\begin{equation}
p(\boldsymbol{x}) = \sum^K_{k=1} \pi_k \mathcal{N}(\boldsymbol{x} \mid \boldsymbol{\mu}_k, \boldsymbol{\Sigma}_k)
\label{eq-marg}
\end{equation}

We can also interpret this as a joint distribution over a space of observable variable $X$ and a space of latent variables $Z$, as follows:
\begin{equation}
  p(\boldsymbol{z}) = \Pi^K_{k=1} \pi^{z_k}_k,
  \label{eq-latent}
\end{equation}
where you should remember that all $z_k$ are zero except for one. The conditional distribution is Gaussian:
\begin{equation}
p(\boldsymbol{x} \mid \boldsymbol{z} ) = \Pi^K_{k=1} \mathcal{N}(\boldsymbol{x} \mid \boldsymbol{\mu}_k, \boldsymbol{\Sigma}_k),
\label{eq-conditional}
\end{equation}
and since:
$$
p(\boldsymbol{x}) = \Pi^K_{k=1} \pi^{z_k}_k
$$
we can write down the joint probability distribution:
\begin{equation}
  p(\boldsymbol{x}, \boldsymbol{z} \mid \boldsymbol{\mu}, \boldsymbol{\Sigma}, \boldsymbol{\pi}) =  \Pi^K_{k=1} \pi^{z_k} 
    \mathcal{N}(\boldsymbol{x} \mid \boldsymbol{\mu}_k, \boldsymbol{\Sigma}_k)^{z_k},
    \label{eq-joint}
\end{equation}
where we have introduced
$$
\boldsymbol{\mu} = (\boldsymbol{\mu}_1, \cdots, \boldsymbol{\mu}_K)^T
$$
and
$$
\boldsymbol{\Sigma} = (\boldsymbol{\Sigma}_1, \cdots \boldsymbol{\Sigma}_K)^T
$$
We are now able to interpret Eq. \ref{eq-marg} as the marginalisation of Eq. \ref{eq-joint}:
$$
p(\boldsymbol{x}) = \sum_{\boldsymbol{z}} p(\boldsymbol{z}) p( \boldsymbol{x} \mid \boldsymbol{z}) = \sum^K_{k=1} \pi_k \mathcal{N}( \boldsymbol{x} \mid \boldsymbol{\mu}_k, \boldsymbol{\Sigma}_k)
$$
This shows that it makes sense to define a latent variable $\boldsymbol{z}_i$ with each observable variable $\boldsymbol{x}_i$ and that we can write down a joint
distribution over both sets of variables. It may not yet be clear why this is useful, given that the GMM Eq. \ref{eq-marg} does not apparently involve them, but
their role will become clear later.

It will also turn out to be useful to introduce $p(z_k = 1 \mid \boldsymbol{x})$. It is sufficiently important to receive its own symbol:
\begin{equation}
\gamma(z_k) \equiv p(z_k = 1 \mid \boldsymbol{x})
\label{eq-respsymb}
\end{equation}
and is called the \emph{responsibility} that component $k$ takes for explaining data point $\boldsymbol{x}$. Intuitively, this is related
to the distance of a data point to the other points of that cluster. This can be made precise through  Bayes' Theorem, as it can be related to
$p(\boldsymbol{x} \mid z_k = 1)$ and therefore:
\begin{align}
  \gamma(z_k)   = p(z_k = 1 \mid \boldsymbol{x} ) & = \frac{p(z_k = 1)p(\boldsymbol{x} \mid z_k = 1) }{\sum^K_{j=1} p(z_j = 1)p(\boldsymbol{x} \mid z_j = 1)} \nonumber \\
  & = \frac{\pi_k \mathcal{N}(\boldsymbol{x} \mid \boldsymbol{\mu}_k, \boldsymbol{\Sigma}_k)}{\sum^K_{j=1} \pi_j \mathcal{N}(\boldsymbol{x} \mid \boldsymbol{\mu}_j, \boldsymbol{\Sigma}_j)}
  \label{eq-resp}
\end{align}
We can interpret $\pi_k$ as a prior probability and $\gamma(z_k)$ as the posterior probability that a  given data point $\boldsymbol{x}_k$ belongs to cluster $k$.
\section{Maximum Likelihood Estimation}
\subsection{The Exponential Family}
In Activity 1.6 you have seen that the likelihood is a product of exponentials but that the log likelihood leads to a sum of quadratic terms. The exponential of
the Gaussian is 'eaten' by the logarithm when considering the log likelihood and reduces to sum of quadratic terms, which is easy to handle. Something similar
happened to the Bernoulli process although we did not make this explicit. We will do that here. Remember that the Bernoulli process is given
by:
$$
p(x \mid \mu)  = \mbox{Bern}(x \mid \mu) = \mu^x(1 - \mu)^{1-x}
$$
This can be written as
$$
p( x \mid \mu) = (1 - \mu) \exp{\ln \frac{\mu}{1-\mu}x}
$$
In this form, it is clear that a data point again is a product of exponentials (and other factors) and the  logarithm reduces to a simple sum. Also, remember
that in Unit 2 we studied the loss function that emerged in logistic regression. Again, this likelihood function was a product of (class conditional) Gaussians.
There is an entire  family of distributions the \emph{exponential family} for which the log likelihood is a much simpler expression because it get rids of the
exponential. Bernoulli, Multinomial and Gaussian distributions are examples of this family and together they span such a large part of statistics that
taking the logarithm becomes second nature.

It is recommended that you read section 2.4 of \cite{bishop2006}, which provides a comprehensive overview of the exponential family.


The key point we want to make here is that even when you start out with a joint distribution that has exponential form that as soon as you marginalise, this
no longer is necessarily the case. We will  perform maximum likelihood estimation (MLE) for Eq. \ref{eq-marg}, which we have seen
can be considered to be a marginalisation over latent variables. We will see that although we will employ the same machinery as in Unit 1, the process is harder
because the sum term in Eq. \ref{eq-marg} prevents the logarithm from cancelling the exponentials.  We will go through parts of the calculation to make this point explicit.
Later, we will show there is an easier way to do this.

Again, like in Unit 1 we will not assess you on being able to perform these calculation yourself. But if you want to get to the point where you are able to access
the machine learning independently, you should try and follow the logic here.



\subsection{Maximum Likelihood Estimation for the Gaussian Mixture Model}
In this section we will perform maximum likelihood estimation of $\boldsymbol{\mu} = (\boldsymbol{\mu}_1, \cdots, \boldsymbol{\mu}_K)^T$,
and $\boldsymbol{\Pi} = (\pi_1, \cdots, \pi_k)^T$, and simply list the result for $(\boldsymbol{\Sigma}_1, \cdots \boldsymbol{\Sigma}_K)^T$.  We will postpone
its calculation until we have found a simpler way.


First we introduce the shorthand notation:
\begin{align}
  \boldsymbol{\mu} &  = (\boldsymbol{\mu}_1, \boldsymbol{\mu}_2, \cdots, \boldsymbol{\mu}_K)^T \nonumber \\
  \boldsymbol{\Sigma} & = (\boldsymbol{\Sigma}_1, \boldsymbol{\Sigma}_2, \cdots, \boldsymbol{\Sigma}_K)
\end{align}

As usual we will represent the data by an $N \times D$ matrix $\boldsymbol{X}$, where $N$ is the number of data points and $D$ is the dimension of the data space, with the
$i$-th row given by $\boldsymbol{x}_i$.
\begin{equation}
\ln p(\boldsymbol{X} \mid \boldsymbol{\Pi}, \boldsymbol{\mu}, \boldsymbol{\Sigma})  = \sum^N_{i=1} \sum^K_{k=1} \ln \{ \pi_k  \mathcal{N}(\boldsymbol{x}_i \mid \boldsymbol{\mu}_k, \boldsymbol{\Sigma}_k) \}
\label{eq-margllh}
\end{equation}

To ram the point home: there is no easy way to evaluate for example:
$$
\ln (2e^x + 5e^{-x})
$$
The innocuous looking summation over $k$ in Eq. \ref{eq-margllh} prevents the logarithm from cancelling the exponentials. We have to estimate parameters on the likelihood
instead of the log likelihood. We want to calculate:
$$
\frac{\partial \ln p(\boldsymbol{X} \mid \boldsymbol{\Pi}, \boldsymbol{\mu}, \boldsymbol{\Sigma})}{\partial \boldsymbol{\mu}_j} =
\sum^N_{i=1} \frac{1}{\sum^K_{k=1} \pi_k \mathcal{N}(\boldsymbol{x}_i \ \boldsymbol{\mu}_k, \boldsymbol{\Sigma}_k)}\frac{\partial \mathcal{N}(\boldsymbol{x}_i \mid \boldsymbol{\mu}_j, \boldsymbol{\Sigma}_j)}{\partial \boldsymbol{\mu}_j}
$$
We are now forced to use the definition of the Gaussian directly:
$$
\frac{\mathcal{N}(\boldsymbol{x}_i \mid \boldsymbol{\mu}, \boldsymbol{\Sigma}) }{ \partial \boldsymbol{\mu}} = (2 \pi)^{-D/2)} \mbox{det}(\Sigma)^{-\frac{1}{2}} \frac{\partial}{\partial \boldsymbol{\mu}} e^{-\frac{1}{2}(\boldsymbol{x}_i - \boldsymbol{\mu})^T \boldsymbol{\Sigma}^{-1} (\boldsymbol{x}_i - \boldsymbol{\mu})}
$$
We have
$$
\frac{\partial}{\partial \boldsymbol{\mu}} e^{-\frac{1}{2}(\boldsymbol{x} - \boldsymbol{\mu})^T \boldsymbol{\Sigma}^{-1}(\boldsymbol{x} - \boldsymbol{\mu})} =
 -\frac{1}{2} e^{-\frac{1}{2}\{\cdots \}}\frac{\partial}{\partial \boldsymbol{\mu}} [ (\boldsymbol{x} - \boldsymbol{\mu})^T \boldsymbol{\Sigma}^{-1} (\boldsymbol{x} - \boldsymbol{\mu}) ]
 $$
 For the last differentiation we can employ the matrix identity that we already used in Unit 1:
 $$
 \frac{\partial}{\partial \boldsymbol{a} } (\boldsymbol{a}^T \boldsymbol{A} \boldsymbol{a}) = (\boldsymbol{A} + \boldsymbol{A}^T)\boldsymbol{a},
 $$
 and using the symmetry of $\boldsymbol{\Sigma} (= \boldsymbol{\Sigma}^T)$, we find that:
 \begin{equation}
   \frac{\partial}{\partial \boldsymbol{\mu} } \mathcal{N}(\boldsymbol{x}_i \mid \boldsymbol{\mu}, \boldsymbol{\Sigma} ) = \mathcal{N}(\boldsymbol{x}_i \mid \boldsymbol{\mu}, \boldsymbol{\Sigma} ) \boldsymbol{\Sigma}^{-1}( \boldsymbol{x}_i - \boldsymbol{\mu})
   \label{eq-diffmu}
 \end{equation}
 Substituting Eq. \ref{eq-diffmu} into \ref{eq-margllh} then gives:
 $$
 \frac{\partial}{\partial \boldsymbol{\mu}_j} \ln p(\boldsymbol{x} \mid \theta) = \sum^N_{i=1} \frac{\pi_j \mathcal{N}(\boldsymbol{x}_i \mid \boldsymbol{\mu}_j, \boldsymbol{\Sigma}_j )}{ \sum^K_{k=1} \pi_k \mathcal{N}(\boldsymbol{x}_i \mid \boldsymbol{\mu}_k, \boldsymbol{\Sigma}_k)} \boldsymbol{\Sigma}^{-1}_j(\boldsymbol{x}_i - \boldsymbol{\mu}_j)
 $$
 Here we are starting to adopt a single symbol, $\theta$, to signify all parameters. This is still a bit unwieldy, so we  introduce
 the symbol:
 \begin{equation}
   \gamma(z_{ij}) \equiv  \frac{\pi_j \mathcal{N}(\boldsymbol{x}_i \mid \boldsymbol{\mu}_j, \boldsymbol{\Sigma}_j) }{\sum^K_{k=1} \pi_k \mathcal{N}(\boldsymbol{x}_i \mid \boldsymbol{\mu}_k, \boldsymbol{\Sigma}_k )},
 \end{equation}
 which happens to be the responsibility for data point $\boldsymbol{x}_i$ (see Eq. \ref{eq-resp}).

The MLE for $\boldsymbol{\mu}_j$  can now be found by:
$$
 \frac{\partial}{\partial \boldsymbol{\mu}_j} \ln p(\boldsymbol{x} \mid \theta) \mid_{\boldsymbol{\mu}_j = \boldsymbol{\mu}_{ML,j}} = 0,
$$
   which leads to:
   $$
   \sum^N_{i=1} \gamma(z_{ij}) \boldsymbol{\Sigma}^{-1}_j (\boldsymbol{x}_i  - \boldsymbol{\mu}_{ML,j}) = 0,
   $$
   which can be easily solved giving:
   \begin{equation}
     \boldsymbol{\mu}_j = \frac{1}{N_j} \sum^N_{i=1} \gamma(z_{ij}) \boldsymbol{x}_i,
     \label{eq-mu}
   \end{equation}
   with:
   $$
   N_j = \sum^N_{i=1} \gamma(z_{ij})
   $$
   A similar exercise can be carried out for the covariance matrix, using similar techniques as in Unit 1,
   but requiring much more effort. We just state the result:
 \begin{equation}
     \boldsymbol{\Sigma}_k = \frac{1}{N_k} \sum^N_{i=1} \gamma( z_{ik} ) (\boldsymbol{x}_i - \boldsymbol{\mu}_k)(\boldsymbol{x}_i - \boldsymbol{\mu}_k)^T,
     \label{eq-sigma}
 \end{equation}
 where the $\boldsymbol{\mu}_k$ are the maximum likelihood estimates given above.

 Finally, the MLE for $\boldsymbol{\pi}$ need to consider the constraint
 $$
 \sum^K_{k=1} \pi_k = 1.
 $$
 This is done by incorporating a Lagrangian multiplier. Again,
 you are not assessed on knowing this technique, this derivation is given for completeness.

 In order to determine $\boldsymbol{\pi}$ we need to maximise:
 $$
 \ln p( \boldsymbol{X} \mid \theta) + \lambda( \sum^K_{k=1} \pi_k -1 )
 $$
 Differentiating this quantity is easy and allows us to eliminate $\lambda$
 $$
 0 = \sum^N_{i=1} \frac{ \mathcal{N}(\boldsymbol{x}_i \mid \boldsymbol{\mu}_k, \boldsymbol{\Sigma}_k)}{ \sum_j \pi_j \mathcal{N}(\boldsymbol{x}_i \mid \boldsymbol{\mu}_j, \boldsymbol{\Sigma}_j )} + \lambda
 $$
 Multiplying both sides by $\pi_k$ and summing over $k$ while using  $\sum^K_{k=1} \pi_k = 1$, one finds
 $$
 \lambda = -N
 $$
 Eliminating $\lambda$ one finds:
 \begin{equation}
   \pi_k = \frac{N_k}{N}
   \label{eq-pi}
 \end{equation}

 Equations \ref{eq-mu}, \ref{eq-sigma} and \ref{eq-pi} do not provide a closed form solution for the parameters. To see this, consider that the responsibilities depend
 $\boldsymbol{\mu}, \boldsymbol{\Sigma}$, and $\boldsymbol{\pi}$, but they in turn are used in the MLE of these parameters. The best we can hope for is an iterative scheme.
 We will present this in the next section.

 \section{E/M Algorithm for Gaussian Mixture Models}
 The Expectation/Maximisation (E/M) algorithm for Gaussian Mixture Models can now be presented as follows:

\begin{enumerate}
\item Determine $K$, the number of clusters. Variants of the E/M algorithm exist that automatically phase out irrelevant clusters  if $K$ is chosen too high but in the
  basic version of the algorithm that we now will discuss a judicious choice for $K$ is important. This may require some experimentation, for example, if the
  clusters are not of great quality and clustering boundaries appear to be arbitrary it make make sense to try and reduce $K$.

\item Initialise the means $\boldsymbol{\mu}_k$ and covariances $\boldsymbol{\Sigma}_k$ and mixing coefficients $\pi_k$. This can be done by picking
  random values, or starting with running the $K$-nearest neighbour algorithm. To focus on essentials, we will assume random initialisation here.

\item {\bf E-step}.  Evaluate the responsibilities:
  $$
  \gamma(z_{ik}) = \frac{ \pi_k \mathcal{N}(\boldsymbol{x}_i \mid \boldsymbol{\mu}_k, \boldsymbol{\Sigma}_k) }{\sum_j \pi_j \mathcal{N}(\boldsymbol{x}_j \mid \boldsymbol{\mu}_j, \boldsymbol{\Sigma}_j)}
  $$
  
  This step evaluates for each data point $\boldsymbol{x}_i$ how much responsibility cluster $k$ assumes for explaining this data point. In actual realisation
  each data point belongs to one and only one cluster, but the higher $\gamma(z_{ik})$, the more likely the data point belongs to cluster $k$. Given that
  this step entails the calculation of $NK$ numbers, this is usually the most time consuming step of the algorithm.

\item {\bf M-step}. Re-estimate the parameters using the responsibilities:

  \begin{align}
    \boldsymbol{\mu}^{new}_k &= \frac{1}{N_k} \sum^N_{i=1} \gamma(z_{ik}) \boldsymbol{x}_i \\
    \boldsymbol{\Sigma}^{new}_k & = \frac{1}{N_k} \sum^N_{i=1} \gamma(z_{ik}) (\boldsymbol{x}_i - \boldsymbol{\mu}^{new}_i)^T(\boldsymbol{x}_i - \boldsymbol{\mu}^{new}_i)\\
    \pi_k &= \frac{N_k}{N}
  \end{align}
  with
  \begin{equation}
    N_k = \sum^N_{i=1} \gamma(z_{ik})
  \end{equation}

\item Evaluate the log likelihood:
  \begin{equation}
    \ln p( \boldsymbol{X} \mid \boldsymbol{\mu}, \boldsymbol{\Sigma}, \boldsymbol{\pi} ) = \sum^N_{i=1} \ln \left\{ \sum^K_{k=1} \pi_k \mathcal{N}(\boldsymbol{x}_i \boldsymbol{\mu}_k,\boldsymbol{\Sigma}_k) \right\}
    \label{eq-mllh}
    \end{equation}
  Compare this to the previous evaluation of the log likelihood for convergence. If not satisfied repeat from step 3 onwards.
    
\end{enumerate}
\subsection{Numerical Stability}
An unlucky choice of initial parameters may cause numerical instability of the algorithm. If  cluster centre $\boldsymbol{\mu}_l$ is chosen that nearly coincides with one
of the data points, the algorithm can maximise the log likelihood by shrinking the variance of cluster $l$, whilst moving $\boldsymbol{\mu}_l$ even closer to the data point.
A detailed explanation  is given in Section 9.2.1 of \cite{bishop2006}. There are various ways to avoid this problem, the easiest is by starting the cluster a number of times
from different random values, something that is prudent anyway. To avoid this particular problem, you could pre-check whether the $\boldsymbol{\mu}_i$ coincide with any data
points. A Bayesian version of this algorithm would not suffer from this problem, which is typical of problems with maximum likelihood estimation, but a full
Bayesian treatment of the algorithm is hard. However, a so-called \emph{variational inference} approach is available. In \emph{Activity 5.B} we will employ
\emph{scikit-learn} and you will find that
a K-means clustering available as a \emph{GaussianMixture} object and the variational version is available as a \emph{BayesianGausianMixture} object. At the expense
of setting more input parameters, which specifies a \emph{prior} on the mixture components and increased computing time, you will have a stabler algorithm, which also
regularises away irrelevant clusters, i.e. if the choice for $K$ is unfortunate and too large, the variational variant will produce some
clusters with very small mixture components that can be ignored in subsequent analysis.

The subject of variational inference is too technical to consider in this unit, but the \emph{Evidence Lower Bound} or \emph{ELBO} plays a central role in this technique. We will
discuss the ELBO below.


\section{GMMs in Practice}
\begin{figure}[!ht]
  \begin{center}
    \includegraphics[width=0.7\textwidth]{start.pdf}
  \end{center}
  \caption{The E/M algorithm in action. At the start random values determine the position of the mean of Gaussian clusters, whereas the covariances have
    been chosen to be isotropic with the same parameter (isotropic means that no direction is preferred, therefore the clusters appear circular). The
    clusters are indicated by isoprobability contour lines.}
  \label{fig-start}
\end{figure}
\begin{figure}[!ht]
  \begin{center}
    \includegraphics[width=0.7\textwidth]{gmmpractice.pdf}
  \end{center}
  \caption{As the algorithm progresses, the position of the Gaussian estimates are adapted. The points of the dataset are coloured according to the values of the responsibility
    in each step, with points mainly associated with cluster 1, being mostly red, cluster 2 green, etc. This colour scheme, which uses RGB values
    that must sum to one works nicely here, because there are three clusters.}
  \label{fig-gmmpractice}
\end{figure}

In Fig. \ref{fig-start}, you can  see the same dataset that was shown in in Fig. \ref{fig-mixtexample}, as well as the initial guess for three Gaussian clusters. Isoprobability
lines are shown for each of these clusters. In Fig. \ref{fig-gmmpractice}, you can see the algorithm in action and observe its performance on each subsequent step. In
Fig. \ref{fig-end}, you can see the end result, the placement of the clusters as well as the responsibility that each cluster takes for each data point. Since we
have three clusters, we can use RGB codes hat add to one in this case. When there are more or less clusters this kind of visualisation does not work. An alternative
could be to use an orientation for each cluster and use saturation level to express the numerical value of the responsibility.
\begin{figure}[!ht]
  \begin{center}
    \includegraphics[width=0.7\textwidth]{end.pdf}
  \end{center}
  \caption{The final result in close up. Note that the clusters can be seen to overlap. This is also visible in the responsibilities, although
    the cluster that is closer usually dominates. Note that the covariant structure is captured well. We started with three isotropic clusters,
  but the algorithm has found covariance matrices that represent the data well.}
  \label{fig-end}
\end{figure}

\section{The E/M Algorithm in General}
\label{sec-gmm}
The E/M algorithm provides a satisfactory solution for the GMM, and as we will see for other mixture models as well. However, at this stage it may not be clear
why the algorithm converges. our treatment has been entirely heuristic. Also, it may not be clear what the advantage is of considering the latent variables
$\boldsymbol{z}$. They feature very indirectly in our estimates of the mixing coefficients $\boldsymbol{\pi}$ but they do not seem to be essential.

We will show in this section that by using latent variables the same problem becomes more tractable. It will lead to the same algorithm, so at this stage
it may still not be entirely clear that latent variables bring a crucial advantage, but in the next section we will show that using this approach
we can proof convergence of the algorithm, and that this proof applies very generally, not just to GMMs. It also sets the stage for a Bayesian
approach to mixture models.

In the previous section, we have seen that the maximum likelihood estimation of a marginalised distribution is hard: we tried to maximise
$$
\ln p(\boldsymbol{X} \mid \theta) = \ln \sum_{\boldsymbol{Z}} p(\boldsymbol{X}, \boldsymbol{Z} )
$$
on the marginalised probability, i.e. with the summation already performed and found that the logarithm no longer works on the distributions directly. It is
generally true that marginalised probabilities, and this includes conditional probabilities, are much harder to handle.

In Unit 4 you will have seen that complex joint probability distributions can be built from networks where the edges represent conditional probabilities. Evaluating these
joint probability distribution functions is cheap, and if necessary, they can be sampled with relative ease. To answer interesting questions, however,
we want to marginalise over most variables. Or we want to answer questions as: 'where is 90 \% of our probability mass located?'. This requires
integration in a multi-dimensional space which is typically analytically intractable and numerically very expensive.

So, whilst $p(x_1, x_2, x_3, x_4, x_5, x_6)$ may be easy to evaluate,
$$
p(x_1, x_3) = \int d x_2 d x_4,  d x_5, d x_6 p(x_1, x_2, x_3, x_4, x_5, x_6) 
$$
is not and neither is:
$$
p(x_2, x_3, x_4 \mid x_1, x_5, x_6) = \frac{p(x_1, x_2, x_3, x_4, x_5, x_6)}{\int d x_1 d x_5 d x_6 p(x_1, x_2, x_3, x_4, x_5, x_6)}  
$$
This is possibly counter intuitive: all knowledge is contained in the joint probability distribution. However, when it comes to extracting
numerical quantities from them, the rule of the thumb is: $\sum_{Z} p(\boldsymbol{X}, \boldsymbol{Z})$: \emph{bad}, $p(\boldsymbol{X}, \boldsymbol{Z})$: \emph{good}.

Let us write down the joint probability distribution for a GMM in terms of both the observable variables $\boldsymbol{X}$, and $\boldsymbol{Z}$:
$$
p(\boldsymbol{X}, \boldsymbol{Z} \mid \boldsymbol{\mu}, \boldsymbol{\Sigma}, \boldsymbol{\pi} ) = \sum^N_{i=1} \pi^{z_{ik}}_k \mathcal{N} (\boldsymbol{x}_i \mid \boldsymbol{\mu}_k, \boldsymbol{\Sigma}_k)^{z_{ik}}
$$

Remember that $\boldsymbol{X} = \left\{ \boldsymbol{x}_1, \boldsymbol{x}_2, \cdots \boldsymbol{x}_N \right\}$ is a set of $N$ points in data space and that
$\boldsymbol{Z} = \left\{ \boldsymbol{z}_1, \cdots, \boldsymbol{z}_N \right\}$ is a set of $N$ $K$-dimensional vectors which have all but one element equal to zero, which indicates to which cluster each data point belongs. Also remember that
we do not have access to these vectors, as they are unobserved.
  

After taking the logarithm:
\begin{equation}
  \ln p(\boldsymbol{X}, \boldsymbol{Z} \mid \boldsymbol{\mu}, \boldsymbol{\Sigma}, \boldsymbol{\pi} ) = \sum^N_{i=1} \sum^K_{k=1}  z_{ik} \left\{ \ln \pi_k + \ln \mathcal{N}(\boldsymbol{x}_i \mid \boldsymbol{\mu}_k, \boldsymbol{\Sigma}_k)\right\}
  \label{eq-fulllh}
\end{equation}
Here $z_{ik}$ is the $k$-th component of $\boldsymbol{z}_i$.

Compared to Eq. \ref{eq-margllh} we see that the summation over $k$ and the logarithm are reversed. The logarithm now directly operates on the exponential.

Suppose now, someone would label the data, i.e. gives us access to the variables $\left\{ \boldsymbol{Z} \right \}$. We could then sort
the terms  in Eq. \ref{eq-fulllh} such that all terms that have $z_{i1} = 1$ (which implies all other components are 0) come first, then
all terms with $z_{i2} =1$ etc., i.e., with access to the labels we would be able to sort the terms in the order of which cluster they appear.
But this would amount to a sum of $K$ independent mixtures. We would be able to perform MLE on each mixture individually, and obtain $K$ estimates for
$\boldsymbol{\mu}_k$ and $\boldsymbol{\Sigma}_k$.

This is not surprising, but the full joint probability distribution shows the relationship between the labelled case, i.e. where we are given which cluster
each data point belongs to and the unlabelled case where the $\boldsymbol{z}_i$  are not observed. In the marginalised distribution, Eq. \ref{eq-margllh}, this
relationship is not visible.

Of course, we have to deal with the case where the $\boldsymbol{z}_i$ are unobserved. We do, however, have estimates for the posterior distribution
of the latent variables through Eq. \ref{eq-resp} and we can consider the expectation of the complete data log likelihood with respect to the posterior distributions of the
latent variables.

To do this requires careful interpretation of the notation, but is not hard. First of all we note that the latent variable distribution is given by Eq. \ref{eq-latent}:
$$
p(\boldsymbol{z}) = \Pi^K_{k=1} \pi^{z_k}_k.
$$
and that the conditional probabilities are Gaussian as per Eq. \ref{eq-conditional}
$$
p(\boldsymbol{x} \mid \boldsymbol{z} ) = \Pi^K_{k=1} \mathcal{N}(\boldsymbol{x} \mid \boldsymbol{\mu}_k, \boldsymbol{\Sigma}_k).
$$
Bayes' rule now gives us an expression for the posterior distribution of latent variables:
\begin{equation}
  p(\boldsymbol{z} \mid \boldsymbol{x} ) \sim p (\boldsymbol{x} \mid \boldsymbol{z}) p( \boldsymbol{z} )
  \label{eq-posterior}
\end{equation}
This becomes a bit abstract; a simple example will help in interpreting this equation. Imagine we have $K=3$, then there are
three different vectors $\boldsymbol{z}_i$, specifically:
\begin{align}
  \boldsymbol{z}_1 & = (1, 0, 0)^T \nonumber \\
  \boldsymbol{z}_2 & = (0, 1, 0)^T \nonumber \\
  \boldsymbol{z}_3 & = (0, 0, 1)^T
  \label{eq-example}
\end{align}
Equation \ref{eq-posterior} states that:
\begin{align}
  p( (1, 0, 0)^T \mid x) & \sim \pi_1 \mathcal{N}(\boldsymbol{x} \mid \boldsymbol{\mu}_1, \boldsymbol{\Sigma}_1) \nonumber \\
  p( (0, 1, 0)^T \mid x) & \sim \pi_2 \mathcal{N}(\boldsymbol{x} \mid \boldsymbol{\mu}_2, \boldsymbol{\Sigma}_2) \nonumber \\
  p( (0, 0, 1)^T \mid x) & \sim \pi_3 \mathcal{N}(\boldsymbol{x} \mid \boldsymbol{\mu}_3, \boldsymbol{\Sigma}_3) \nonumber 
\end{align}
All we have to do to find the actual probabilities is to normalise:
\begin{align}
  p( (1, 0, 0)^T \mid x) & = \frac{ \pi_1 \mathcal{N}(\boldsymbol{x} \mid \boldsymbol{\mu}_1, \boldsymbol{\Sigma}_1)}{\pi_1 \mathcal{N}(\boldsymbol{x} \mid \boldsymbol{\mu}_1, \boldsymbol{\Sigma}_1) + \pi_2 \mathcal{N}(\boldsymbol{x} \mid \boldsymbol{\mu}_2, \boldsymbol{\Sigma}_2) + \pi_3 \mathcal{N}(\boldsymbol{x} \mid \boldsymbol{\mu}_3, \boldsymbol{\Sigma}_3)} \nonumber \\
  p( (0, 1, 0)^T \mid x) & = \frac{ \pi_2 \mathcal{N}(\boldsymbol{x} \mid \boldsymbol{\mu}_2, \boldsymbol{\Sigma}_2)}{\pi_1 \mathcal{N}(\boldsymbol{x} \mid \boldsymbol{\mu}_1, \boldsymbol{\Sigma}_1) + \pi_2 \mathcal{N}(\boldsymbol{x} \mid \boldsymbol{\mu}_2, \boldsymbol{\Sigma}_2) + \pi_3 \mathcal{N}(\boldsymbol{x} \mid \boldsymbol{\mu}_3, \boldsymbol{\Sigma}_3)} \nonumber \\
  p( (0, 0, 1)^T \mid x) & = \frac{ \pi_3 \mathcal{N}(\boldsymbol{x} \mid \boldsymbol{\mu}_3, \boldsymbol{\Sigma}_3)}{\pi_1 \mathcal{N}(\boldsymbol{x} \mid \boldsymbol{\mu}_1, \boldsymbol{\Sigma}_1) + \pi_2 \mathcal{N}(\boldsymbol{x} \mid \boldsymbol{\mu}_2, \boldsymbol{\Sigma}_2) + \pi_3 \mathcal{N}(\boldsymbol{x} \mid \boldsymbol{\mu}_3, \boldsymbol{\Sigma}_3)}
\label{eq-post}
\end{align}

It is easy to see that in the general case
\begin{equation}
p(\boldsymbol{z} \mid \boldsymbol{x} ) = \frac{ \Pi^K_{k=1}( \pi_k \mathcal{N}( x \mid \boldsymbol{\mu}_k, \boldsymbol{\Sigma}_k))^{z_{k}}}{\sum^K_{j=1} \pi_j \mathcal{N}(\boldsymbol{x} \mid \boldsymbol{\mu}_j, \boldsymbol{\Sigma}_j)}
\end{equation}
Here the $z_k$ are the components of the vector $\boldsymbol{z}$.
The result should make sense. What Eq. \ref{eq-post} expresses is that it is more likely to belong to a particular component if you're closer to it and further
removed for the others.


We are almost set. We will need one more building block and that is the expectation value of the vector $\boldsymbol{z}$ with respect to the posterior latent
distribution. Again, an example will clarify that this is a very simple concept. By definition the expectation  value is given by:
\begin{align}
\mathbb{E}_{\boldsymbol{z}}[\boldsymbol{z}] &= 
\left( \begin{array}{c} 1 \\ 0 \\ 0 \end{array} \right)p(\left( \begin{array}{c} 1 \\ 0 \\ 0 \end{array} \right)  \mid \boldsymbol{x} ) +  \left( \begin{array}{c} 0 \\ 1 \\ 0 \end{array} \right)p(\left( \begin{array}{c} 0 \\ 1 \\ 0 \end{array} \right) \mid \boldsymbol{x} ) +  \left( \begin{array}{c} 0 \\ 0 \\ 1 \end{array} \right)p(\left( \begin{array}{c} 0 \\ 0 \\ 1 \end{array} \right)  \mid \boldsymbol{x} ) \nonumber \\
& = \left(\begin{array}{c} \frac{\pi_1 \mathcal{N}(\boldsymbol{x}\mid \boldsymbol{\mu}_1, \boldsymbol{\Sigma}_1)}{\sum^K_{j=1} \pi_j \mathcal{N}(\boldsymbol{x} \mid \boldsymbol{\mu}_j, \boldsymbol{\Sigma}_j)} \\ \frac{\pi_2 \mathcal{N}(\boldsymbol{x} \mid \boldsymbol{\mu}_2, \boldsymbol{\Sigma}_2)}{\sum^K_{j=1} \pi_j \mathcal{N}(\boldsymbol{x} \mid \boldsymbol{\mu}_j, \boldsymbol{\Sigma}_j)} \\ \frac{\pi_3\mathcal{N}(\boldsymbol{x} \mid \boldsymbol{\mu}_3, \boldsymbol{\Sigma}_3)}{\sum^K_{j=1} \pi_j \mathcal{N}(\boldsymbol{x} \mid \boldsymbol{\mu}_j, \boldsymbol{\Sigma}_j)} \end{array} \right)
\end{align}

Note that the components of this expectation are nothing but the responsibilities defined in Eq. \ref{eq-resp},
in other words:
\begin{equation}
\mathbb{E}_{\boldsymbol{Z}}[\boldsymbol{z}]_k = \gamma(z_k)
\label{eq-gamma}
\end{equation}

So far, this may look like wrapping relatively obvious facts in highly abstract notation, but these exercises allow us to evaluate the expectation value of the log likelihood
with respect to the posterior latent variable distribution in a straightforward way. Consider the full data log likelihood Eq. \ref{eq-fulllh} and taking the expectation value,
the result is simple:
\begin{equation}
  \mathbb{E}_{\boldsymbol{Z}}[\ln p( \boldsymbol{X}, \boldsymbol{Z} \mid \boldsymbol{\mu}, \boldsymbol{\Sigma}, \boldsymbol{\pi})] = \sum^N_{i=1} \sum^K_{k=1} \gamma( z_{ik} )
  \left\{ \ln \pi_k  + (\boldsymbol{x}_i - \ln \mathcal{N}( \boldsymbol{x}_i \mid \boldsymbol{\mu}_k, \boldsymbol{\Sigma}_k) \right\}
  \label{eq-zlogllh}
\end{equation}
Note how our work leading up to Eq. \ref{eq-gamma} pays off: the $z_{ik}$ in Eq. \ref{eq-fulllh} are the only dependants on $\boldsymbol{z}$.


In \emph{Activity 5.c}, you will be asked to perform maximum likelihood estimation with the aid of Eq. \ref{eq-zlogllh}. You will note that this is much simpler and
that you can reuse the results of Unit 1. This is an important example of how joint distributions are easier to handle than marginal ones.

The approach in this section shows why the algorithm is called the Expectation/Maximisation algorithm. First, the \emph{expectation} of the joint probability
is calculated. The likelihood of the resulting distribution is then \emph{maximised}. The two step approach that was already recognisable in the K-means
clustering is clearly present, but now better motivated. In Sec. \ref{sec-elbo}, we will show that this approach always leads to an increased likelihood of the model.
So far, this is not obvious. In the M-step, the likelihood will increase by definition, but the E-step, taking the expectation with respect to the posterior latent variables
might undo the previous optimisation step in theory. We will see this is not the case. The manner of how this is shown is very important to modern
machine learning.

Finally, we mention in passing that to calculate that the likelihood of the posterior distribution of latent variables is simply the product

\section{The Evidence Lower Bound (ELBO)}
\label{sec-elbo}
The argument we will present here is abstract but simple. The main problem in latent variant models such as the GMM is to estimate the distribution of the
latent variables $p(\boldsymbol{Z} \mid \boldsymbol{X}, \theta)$, as this is the one we cannot observe. We usually are forced to assume a distribution
$q(\boldsymbol{Z})$, which \emph{hope} will be close to the true distribution of latent variables. We will show that for any such $q(\boldsymbol{Z})$
the following decomposition is valid:
\begin{equation}
  \ln p(X \mid \theta ) = \sum_{\boldsymbol{Z}} q(\boldsymbol{Z}) \ln \left\{ \frac{p(\boldsymbol{X},\boldsymbol{Z}) }{q (\boldsymbol{Z}) }  \right\}
  - \sum_{\boldsymbol{Z}} q(\boldsymbol{Z} ) \ln \left\{ \frac{p(\boldsymbol{Z} \mid \boldsymbol{X}, \theta)}{q(\boldsymbol{Z})} \right\}
  \label{eq-elbo}
\end{equation}
We will first show that this is true. From Unit 1, we remember the product rule:
$$
p(\boldsymbol{X} ,\boldsymbol{Z} \mid \theta)  = p(\boldsymbol{Z}\mid \boldsymbol{X}, \theta) p(\boldsymbol{X} \mid \theta)
$$
and therefore, taking the logarithm:
$$
\ln p(\boldsymbol{X} ,\boldsymbol{Z} \mid \theta)  = \ln p(\boldsymbol{Z}\mid \boldsymbol{X}, \theta) + \ln p(\boldsymbol{X} \mid \theta)
$$
We can expand the first term on the right hand side of Eq. \ref{eq-elbo}
\begin{align}
\sum_{\boldsymbol{Z}} q(\boldsymbol{Z}) \ln p(\boldsymbol{X}, \boldsymbol{Z} \mid \theta) - \sum_{\boldsymbol{Z}} q(\boldsymbol{Z}) \ln q(\boldsymbol{Z})  &= \nonumber \\
\sum_{\boldsymbol{Z}} q(\boldsymbol{Z}) \{ \ln p(\boldsymbol{Z}\mid \boldsymbol{X}, \theta) + \ln p(\boldsymbol{X} \mid \theta) \} - \sum_{\boldsymbol{Z}} q(\boldsymbol{Z}) \ln q(\boldsymbol{Z}
)
\end{align}
Adding the second term on the right hand side of Eq. \ref{eq-elbo} cancels most terms. The only remaining contribution is:
$$
\sum_{\boldsymbol{Z}} q(\boldsymbol{Z})  \ln p(\boldsymbol{X} \mid \theta) 
$$
But $\sum_{\boldsymbol{Z}} q(\boldsymbol{Z}) =1$, because it is a probability distribution and $\boldsymbol{Z}$ does not feature anywhere else.
So Eq. \ref{eq-elbo} is true.

We recall the definition of the Kullback-Leibler divergence from Unit 1:
$$
\mbox{KL}( q \mid \mid p) = -\sum_{\boldsymbol{Z}} q(\boldsymbol{Z}) \ln \left\{ \frac{p(\boldsymbol{Z} \mid \boldsymbol{X}, \theta)}{q(\boldsymbol{Z})} \right\}
$$
and introduce the so-called \emph{evidence lower bound}:
$$
\mathcal{L}(q, \theta) = \sum_{\boldsymbol{Z}} q(\boldsymbol{Z}) \left\{ \frac{p(\boldsymbol{X},\boldsymbol{Z} \mid \theta) }{q(\boldsymbol{Z})} \right\}
$$
so that we can write:
\begin{equation}
  \ln p(\boldsymbol{X} \mid \theta) = \mathcal{L}(q, \theta) + \mbox{KL}(q\mid\mid p)
  \label{eq-decom}
 \end{equation}
Despite the abstract notation, some of the actors should look familiar. The goal is to maximise the log likelihood $p(\boldsymbol{X} \mid \theta)$. With the
GMM, we have given an example for why this is difficult to do directly and for more complex cases this is generally intractable. In the lower bound, we recognise the
expectation of the joint distribution with respect to an estimate for the latent variable distribution $q( \boldsymbol{Z})$, which we found much easier
to optimise. The second term on the right hand side, $\mbox{KL}(q \mid \mid p)$ is a measure for how much $q$ differs from $p$. Moreover, this measure
is only zero when $q = p$ and otherwise positive.

This justifies the name lower bound for $\mathcal{L}$: whatever value of the likelihood $p(\boldsymbol{X} \mid \theta )$ is, it is always higher than the lower bound as the
KL term gives adds positively to the lower bound to $\mathcal{L}$, unless $q = p$, which only happens if estimate $q(\boldsymbol{Z})$
is equal to the true posterior distribution of the latent variable: $p(\boldsymbol{Z} \mid \boldsymbol{X}, \theta)$.

We can now identify the E and M steps in the E/M algorithm as follows. Suppose we have an estimate for the parameters $\theta^{old}$, which
results in $\ln p(\boldsymbol{X} \mid \theta^{old})$ for the value of the log likelihood of the marginal distribution. The E-step consists in adapting
$q(\boldsymbol{Z})$ so that it matches the posterior likelihood $p(\boldsymbol{Z} \mid \boldsymbol{X})$. In the GMM version of E/M this is done by
updating the responsibilities. As explained in Section \ref{sec-gmm}, updating the responsibilities and calculating the posterior likelihood are equivalent. At this
stage $q=p$, and the KL-divergence is equal to 0, \emph{so for this choice of $q$, for parameters $\theta^{old}$ we have}


\begin{equation}
\ln p(\boldsymbol{X} \mid \theta^{old}) = \sum_{\boldsymbol{Z}} q(\boldsymbol{Z})  \ln \frac{p(X,Z \mid \theta^{old})}{q(Z)}
\label{eq-temp}
\end{equation}
We now find parameters $\theta^{new}$ that maximise $\sum_{\boldsymbol{Z}} p( \boldsymbol{X}, \boldsymbol{Z} \mid \theta)$. Because of the equality in Eq. \ref{eq-temp},
if we manage to find parameters $\theta^{new}$ that maximise the lower bound, i.e. increase it (if we cannot find parameters that increase the lower bound we should stop),
then necessarily for our current value of $q(\boldsymbol{Z})$, they also increase the marginal likelihood.

It is worthwhile making that explicit. Consider the situation immediately after an E-step.
At this stage $q(\boldsymbol{Z}) = p(\boldsymbol{Z} \mid \boldsymbol{X}, \theta_{old})$.
Substituting this in the lower bound, we can write:
\begin{align}
  \mathcal{L}(q, \theta) &  = \sum_{Z}p(\boldsymbol{Z} \mid \boldsymbol{X}, \theta^{old}) \ln p (X, Z \mid \theta) - \sum_{\boldsymbol{Z}} p (\boldsymbol{Z} \mid \boldsymbol{X}, \theta^{old})  \ln p( \boldsymbol{Z} \mid \boldsymbol{X}, \theta^{old}) \nonumber \\
  & = \sum_{\boldsymbol{Z}} p( \boldsymbol{Z} \mid \boldsymbol{X}, \theta^{old}) \ln p(\boldsymbol{X}, \boldsymbol{Z} \mid \theta) + \mbox{const}
  \label{eq-dependency}
\end{align}
where the constant is the negative entropy of $q$, which does not depend on $\theta$. Since the $\mbox{KL}$-term is also not dependent on $\theta$, we see that
the depency of $\ln p(\boldsymbol{X} \mid \theta)$ is given by that of $\ln p(\boldsymbol{X}, \boldsymbol{Z} \mid \theta)$.



{\bf After updating the parameters to $\theta^{new}$, the equality
  of Eq. \ref{eq-temp} no longer holds!} For our new parameters, the old responsibilities, which were calculated using $\theta^{old}$ no longer represent the posterior likelihood
$p$ and the KL divergence of $p$ relative to $q$ is greater than 0. We have to calculate new responsibilities, using $\theta^{new}$. We repeat these steps
until in the M-step we are no longer capable of improving the lower bound.

Note that in this process the lower bound functions as a ratchet. Upon calculating the responsibilities, we raise the lower bound so that it equals the marginal probability
as per Eq. \ref{eq-temp}. If we then manage to increase the expectation value of the joint distribution by picking new parameters we raise the likelihood
of the marginal likelihood as well. Then, new responsibilities can be calculated, raising the lower bound to the new marginal likelihood, etc.

The entire process is shown in Fig. \ref{fig-em}. In the E-step, we choose responsibilities such they match the posterior distribution. The lower bound is equal
to $\ln p( \boldsymbol{X} \mid \theta^{old})$, moreover the $\theta$ dependency for it and the lower bound are the same, so the curves do not just intersect, but touch.
In the M-step the lower bound is optimised for $\theta$ and a new value $\theta^{new}$ is found. The lower bound is shown as a quadratic function,
which indeed it is for a mixture of Gaussians. \emph{The discussion of this Section applies to all members of the exponential family, however.}

\begin{figure}[!ht]
  \begin{center}
    \includegraphics[width=0.7\textwidth]{bishop_em_steps.pdf}
  \end{center}
  \caption{This is Figure 9.14 from Bishop (2006). }
  \label{fig-em}
\end{figure}

\section{Variational Inference}
It is possible to give a Bayesian treatment of mixture models. This requires
a prior distribution on the mixing coefficients $\pi$. An advantage of the Bayesian
treatment is that mixture components that are not supported by the data are
regularised away, in much the same way as large weights are regularised away in Bayesian
Linear Regression, at least when they are not supported by subsequent data.
This gives the possibility of estimating the number of components from the data.
After convergence you can find the mean of the mixing coefficient and those
components who have clearly smaller means than the others can be cut away.

The theoretical treatment of variational inference is highly technical and is explained
in Chapter 10 of \cite{bishop2006}. It is well out of scope of this unit, but the
result is an algorithm that is not so different from the one we have discussed so far.
In \emph{Activity 5.B}, we will give a practical example based on a \emph{scikit-learn}
implementation.

    
  
  


  
